% !TEX program = xelatex
\documentclass[12pt]{article}

% الهوامش
\usepackage{geometry}
\geometry{margin=0.8in}

% تنسيقات عامة
\usepackage{xcolor}
\usepackage{fancyhdr}
\usepackage{setspace}
\usepackage[inline]{enumitem}    % enumerate* لسطر واحد
\usepackage{tabularx}

% دعم لغات/اتجاهات
\usepackage{fontspec}
\usepackage{polyglossia}
\usepackage{bidi}                % LTR/RTL

% رياضيات
\usepackage{amsmath, amssymb, mathtools}

% صور
\usepackage{graphicx}

% ضبط اللغات والخطوط
\setmainlanguage[numerals=western]{arabic}   % افتراضي: عربي (أرقام 0-9 إنكليزية)
\setotherlanguage{english}
\newfontfamily\arabicfont[Script=Arabic,Scale=1.05]{Amiri}
\newfontfamily\englishfont{Times New Roman}

% ترويسة/تذييل
\pagestyle{fancy}
\fancyhf{}
\renewcommand{\headrulewidth}{0pt}
\cfoot{\EN{Page } \thepage}

% ===== أدوات اتجاه آمنة =====
% نص إنجليزي باتجاه LTR بشكل آمن وسهل
\newcommand{\EN}[1]{{\begin{LTR}\textenglish{#1}\end{LTR}}}
% نص عربي صريح (إن احتجت لفرضه داخل سياق LTR)
\newcommand{\AR}[1]{{\textarabic{#1}}}

% ===== ماكرو عنوان السؤال (يسار، LTR دائماً) =====
\renewcommand{\qtitle}[1]{%
  \par\noindent{\begin{LTR}\textenglish{\bfseries #1.Question #1}\end{LTR}}\par\vspace{0.4em}%
}

% توافقاً مع مخرجاتك السابقة
\newcommand{\qheader}[1]{\qtitle{#1}}

% ===== عنصر خيار مع دعم مزج اللغات =====
\newcommand{\AutoItem}[1]{#1}

% ===== صفّ الخيارات أفقيًا مع لف تلقائي =====
\newcommand{\renderchoices}[5]{%
  {\begin{LTR}\begin{enumerate*}[
    label=\textbf{(\Alph*)},     % (A) (B) (C) ...
    itemjoin=\hspace{2em},
    itemsep=0.25em,
    leftmargin=0pt
  ]
    \ifx\relax#1\relax\else\item \AutoItem{#1}\fi
    \ifx\relax#2\relax\else\item \AutoItem{#2}\fi
    \ifx\relax#3\relax\else\item \AutoItem{#3}\fi
    \ifx\relax#4\relax\else\item \AutoItem{#4}\fi
    \ifx\relax#5\relax\else\item \AutoItem{#5}\fi
  \end{enumerate*}\end{LTR}}\par
}


% مسافة أسطر مريحة
\setstretch{1.25}

\begin{document}

% ===== Header (مطابق تقريبًا) =====
\noindent
\begin{tabularx}{\textwidth}{@{\extracolsep{\fill}} l c r}
\EN{\textbf{Date:} September 4, 2025} &
{\Large \textbf{\EN{Algorithms -1-}}} &
\EN{University Of Aleppo} \\
\EN{\textbf{Student Name:}} &
\EN{\textbf{Model:} A} &
\EN{\textbf{Duration:} 90 Minutes} \\
\EN{\textbf{Student Number:}} & &
\end{tabularx}

\vspace{1.0cm}

% عنوان قسم الأسئلة LTR
{\EN{\textbf{Select the correct answer for each question:}}}

% ===== أمثلة أسئلة (نفس طابع أمثلتك) =====

\noindent\qheader{1}
{\EN{Calculate the derivative of }} \( f(x) = x^2 \) {\EN{ then find the limit to infinity of the function }} \( \lim_{x \to \infty} f(x) \).
\renderchoices{\EN{14}}{\EN{15}}{\underline{\EN{16}}}{\EN{17}}{\EN{17}}
{\begin{LTR}\hfill\EN{[1 Point]}\end{LTR}}\vspace{0.5em}

\noindent\qheader{2}
{\EN{Please calculate }} \( \int f(x)\,dx \).
\renderchoices{\EN{13}}{\underline{\EN{14}}}{\EN{15}}{\EN{16}}{}
{\begin{LTR}\hfill\EN{[2 Points]}\end{LTR}}\vspace{0.5em}

\noindent\qheader{3}
\EN{Test}
\renderchoices{\EN{13}}{\underline{\EN{14}}}{\EN{15}}{\EN{16}}{\EN{17}}
{\begin{LTR}\hfill\EN{[3 Points]}\end{LTR}}\vspace{0.5em}

\noindent\qheader{4}
\EN{Test2}
\renderchoices{\underline{\EN{1}}}{\EN{2}}{\EN{3}}{\EN{4}}{\EN{5}}
{\begin{LTR}\hfill\EN{[3 Points]}\end{LTR}}\vspace{0.5em}

\noindent\qheader{5}
\EN{Test3}
\renderchoices{\EN{1}}{\EN{2}}{\underline{\EN{3}}}{\EN{4}}{\EN{5}}
{\begin{LTR}\hfill\EN{[1 Point]}\end{LTR}}\vspace{0.5em}

\noindent\qheader{6}
\EN{test4}
\renderchoices{\underline{\EN{1}}}{\EN{2}}{\EN{3}}{\EN{4}}{\EN{5}}
{\begin{LTR}\hfill\EN{[1 Point]}\end{LTR}}\vspace{0.5em}

\noindent\qheader{7}
\EN{Hello!}
\renderchoices{\underline{\EN{1}}}{\EN{2}}{\EN{3}}{\EN{4}}{}
{\begin{LTR}\hfill\EN{[2 Points]}\end{LTR}}\vspace{0.5em}

\noindent\qheader{8}
\EN{Calculate } \(\sum\) \EN{ i}
\renderchoices{\underline{\(n\,(n+1)/3\)}}{\(n\,(n-1)/2\)}{\EN{0}}{\EN{0}}{}
{\begin{LTR}\hfill\EN{[1 Point]}\end{LTR}}\vspace{0.5em}

\noindent\qheader{9}
\EN{Find the area of a rectangle of side } \(a = 6\)
\renderchoices{\EN{34}}{\underline{\EN{36}}}{\EN{40}}{\EN{31}}{}
{\begin{LTR}\hfill\EN{[1 Point]}\end{LTR}}\vspace{0.5em}

\noindent\qheader{10}
\EN{For any positive integer } \(x\) \EN{ what is the value of } \(\iiint x^5\)\EN{?}
\renderchoices{\underline{\EN{30}}}{\EN{40}}{\EN{50}}{\EN{60}}{}
{\begin{LTR}\hfill\EN{[1 Point]}\end{LTR}}\vspace{0.5em}

\noindent\qheader{11}
\EN{Solve for } \(x\) \EN{ in the equation } \(2x + 6 = 0\)
\renderchoices{\EN{-1}}{\EN{-2}}{\underline{\EN{-3}}}{\EN{3}}{}
{\begin{LTR}\hfill\EN{[1 Point]}\end{LTR}}\vspace{0.5em}

\noindent\qheader{12}
\(x=3+6\) \EN{ what is the value of x?}
\renderchoices{\EN{7}}{\underline{\EN{9}}}{\EN{8}}{\EN{10}}{\EN{0}}
{\begin{LTR}\hfill\EN{[1 Point]}\end{LTR}}\vspace{0.5em}

\noindent\qheader{13}
\EN{What is the derivative of } \(x^2\)?
\renderchoices{\(x\)}{\underline{\(2x\)}}{\EN{0}}{\EN{0}}{}
{\begin{LTR}\hfill\EN{[5 Points]}\end{LTR}}\vspace{0.5em}

\end{document}
